\section{Discussion}\label{sec:discussion}

In the family \eqref{eq:fixed}, changing the external circle exchanges
an Euler-characteristic grading with a rotation of a plane. The commutator
deformation accounts for the disappearance of positive excess in the
plane description. The geometry that remains is the Hilbert scheme of
points, and its tangent representation supplies the hook lengths in
the local topological-string partition function. The determinant on the
family of doubling directions distinguishes the balanced splitting:
its degree is zero precisely when the two original normal lines have
equal degree. The first variation in section~\ref{sec:variation}
illustrates what a family retains even when a numerical trace vanishes.

An index on a space of cycles and the Hilbert space of a charged
five-dimensional particle answer different questions. In degree two the
index calculation identifies the locus supporting the connected term
and fixes all its coefficients. A microscopic membrane calculation would
have to determine the spectrum at the threshold where two branes on $C$
can separate. To identify the coefficient at $q^{-2}$ with a particle
spin trace, one would have to establish a normalizable state in the full
unframed theory with the two-particle continuum open, and then derive
its spin from the wrapped-brane dynamics. The value $-4$ computed here
provides an independent comparison for that calculation.
